\documentclass[11pt,a4paper]{article}
\usepackage[utf8]{inputenc}
\usepackage[T1]{fontenc}
\usepackage[catalan]{babel}
\usepackage{lmodern}
\usepackage[sfdefault,lining]{FiraSans}
\usepackage[a4paper,top=2.3cm,bottom=2.2cm,left=1.8cm,right=1.8cm,headheight=1.3cm,headsep=0.4cm,footskip=1cm]{geometry}
\usepackage{amsmath,amssymb}
\usepackage{icomma}
\usepackage{graphicx}
\usepackage[table]{xcolor}
\usepackage{tikz}
\usepackage{fontawesome5}
\usepackage{ccicons}
\usepackage{fancyhdr}
\usepackage{lastpage}
\usepackage{tasks}
\usepackage{needspace}
\usepackage[most]{tcolorbox}
\usepackage{array}
\usepackage[hidelinks,pdftitle={Solucions de la fitxa de la unitat 8: Àrees i perímetres. 1r d'ESO},pdfsubject={Matemàtiques 1r ESO},pdfkeywords={perímetre, àrea, unitats de longitud, unitats de superfície, hectàrea, quadrat, rectangle, romboide, triangle, rombe, trapezi, polígons regulars, apotema, circumferència, nombre pi, cercle, sector circular, corona circular, figures compostes, teorema de Pitàgores, catets, hipotenusa}]{hyperref}

\definecolor{navy}{HTML}{1B2748}
\definecolor{accent}{HTML}{E8603A}
\definecolor{teal}{HTML}{1F9E90}
\definecolor{soft}{HTML}{EEF1F8}
\definecolor{muted}{HTML}{5D6685}

\setlength{\parindent}{0pt}
\setlength{\parskip}{3pt}
\settasks{label=\alph*), label-format=\bfseries\color{navy}, label-width=1.3em, item-indent=1.8em, column-sep=0.6em, after-item-skip=5pt, before-skip=3pt, after-skip=4pt}

% ---------- capçalera i peu ----------
\newcommand{\llicencia}{\href{https://creativecommons.org/licenses/by-nc-sa/4.0/deed.ca}{\ccbyncsa}\ \ Llicència Creative Commons Reconeixement-NoComercial-CompartirIgual 4.0 Internacional (CC BY-NC-SA 4.0)}
\pagestyle{fancy}
\fancyhf{}
\fancyhead[L]{\raisebox{-0.35\height}{\includegraphics[height=0.95cm]{logo-gva.png}}}
\fancyhead[C]{\small\color{muted}Solucions de la fitxa · Unitat 8 · 1r d'ESO}
\fancyhead[R]{\raisebox{-0.35\height}{\includegraphics[height=0.95cm]{logo-vermellar.png}}}
\fancyfoot[L]{\scriptsize\color{muted}\llicencia}
\fancyfoot[R]{\small\color{muted}\thepage}
\renewcommand{\headrulewidth}{0.4pt}
\renewcommand{\headrule}{\color{muted!40}\hrule height \headrulewidth}
\fancypagestyle{primera}{\fancyhf{}\renewcommand{\headrulewidth}{0pt}
  \fancyfoot[L]{\scriptsize\color{muted}\llicencia}
  \fancyfoot[R]{\small\color{muted}\thepage}}

% ---------- nivells i exercicis ----------
\newcommand{\estrelles}[1]{{\color{accent}\ifcase#1\or\faStar\faStar[regular]\faStar[regular]\or\faStar\faStar\faStar[regular]\or\faStar\faStar\faStar\fi}}
\newcounter{exer}
\newcommand{\ex}[1]{\par\needspace{4\baselineskip}\medskip\stepcounter{exer}\noindent{\color{navy}\bfseries\theexer.}\ \ {\small\estrelles{#1}}\ \ }

\newcommand{\seccio}[2]{\par\needspace{10\baselineskip}\bigskip
  \begin{tcolorbox}[colback=navy,colframe=navy,coltext=white,arc=2mm,boxrule=0pt,left=3mm,right=3mm,top=1.2mm,bottom=1.2mm]
  {\large\bfseries #1\quad #2}
  \end{tcolorbox}\vspace{-2pt}}

\newcolumntype{C}[1]{>{\centering\arraybackslash}p{#1}}

% quadre de recordatori (només a la fitxa d'exercicis)
\newcommand{\recorda}[1]{\par\smallskip\begin{tcolorbox}[colback=accent!7,colframe=accent!55,arc=2mm,boxrule=0.5pt,left=2.5mm,right=2.5mm,top=1mm,bottom=1mm]\small{\color{accent}\faLightbulb[regular]}\ \textbf{Recorda:}\ #1\end{tcolorbox}}
\newcommand{\res}[1]{\boxed{#1}}\newcommand{\rest}[1]{\fbox{#1}}

\def\fmneed{14}
% =====================================================================
% Capçalera FentMates de les fitxes: un color per a cada curs.
% Els colors són els de la pàgina inicial (LEVELS a assets/js/materials.js);
% el logo ve de _eines/logo/pdf/ (python3 _eines/logo/genera.py).
%
% Ús, al final del preàmbul de la unitat:
%   % =====================================================================
% Capçalera FentMates de les fitxes: un color per a cada curs.
% Els colors són els de la pàgina inicial (LEVELS a assets/js/materials.js);
% el logo ve de _eines/logo/pdf/ (python3 _eines/logo/genera.py).
%
% Ús, al final del preàmbul de la unitat:
%   % =====================================================================
% Capçalera FentMates de les fitxes: un color per a cada curs.
% Els colors són els de la pàgina inicial (LEVELS a assets/js/materials.js);
% el logo ve de _eines/logo/pdf/ (python3 _eines/logo/genera.py).
%
% Ús, al final del preàmbul de la unitat:
%   \input{../comu/fentmates.tex}
%   \fmcurs{1eso}
% i, al començament del document:
%   \capcalera{Matemàtiques}{1r ESO}{Fitxa d'exercicis · Unitat 2}{Divisibilitat}
% Opcional abans de \input: \def\fmneed{14} (espai mínim abans de cada secció).
% =====================================================================
\makeatletter
\def\fm@logodir{../../logo/pdf}
\providecommand{\fmneed}{10}
\def\fm@cursnom{}

% clau del curs, color i color del text damunt del color
\newcommand{\fm@curs}[3]{\@namedef{fm@color@#1}{#2}\@namedef{fm@text@#1}{#3}}
\fm@curs{1eso}{E6B84A}{navy}
\fm@curs{2eso}{3B82F6}{white}
\fm@curs{3eso}{FF7A50}{white}
\fm@curs{4eso-a}{A78BFA}{white}
\fm@curs{4eso-b}{6D28D9}{white}
\fm@curs{1batx-mi}{9B1C1C}{white}
\fm@curs{1batx-ccss}{EF6B6B}{white}
\fm@curs{2batx-mii}{1E6B47}{white}
\fm@curs{2batx-ccss}{2BB3A3}{white}

\newcommand{\fmcurs}[1]{%
  \@ifundefined{fm@color@#1}{\PackageError{fentmates}{Curs desconegut: #1}{}}{}%
  \edef\fm@tmp{\noexpand\definecolor{curs}{HTML}{\@nameuse{fm@color@#1}}%
               \noexpand\colorlet{curstext}{\@nameuse{fm@text@#1}}}\fm@tmp
  \colorlet{cursfosc}{curs!72!black}%
  \def\fm@logo{\fm@logodir/vertical-#1.pdf}%
  \def\fm@marca{\fm@logodir/marca-#1.pdf}}

% ---------- capçalera de la primera pàgina ----------
% #1 matèria (petit), #2 curs, #3 tipus de document i unitat, #4 títol
\newcommand{\capcalera}[4]{%
  \gdef\fm@cursnom{#2}%
  \par\noindent
  \begin{tcolorbox}[colback=curs,colframe=curs,boxrule=0pt,arc=2.2mm,width=3.5cm,height=2.15cm,
      valign=center,left=3mm,right=2mm,top=1mm,bottom=1mm,nobeforeafter,box align=center]
    \color{curstext}{\footnotesize #1}\par\vspace{1pt}{\fontsize{19}{21}\selectfont\bfseries #2\par}%
  \end{tcolorbox}\hspace{5mm}%
  \begin{minipage}[c]{\dimexpr\linewidth-3.5cm-5mm-2.9cm\relax}
    {\color{cursfosc}\small\bfseries\MakeUppercase{#3}}\par\vspace{2pt}
    {\color{navy}\fontsize{20}{23}\selectfont\bfseries #4\par}
  \end{minipage}\hfill
  \begin{minipage}[c]{2.6cm}\raggedleft\includegraphics[width=2.6cm]{\fm@logo}\end{minipage}\par
  \vspace{2pt}}

% ---------- barres de secció del color del curs ----------
\renewcommand{\seccio}[2]{\par\needspace{\fmneed\baselineskip}\bigskip
  \begin{tcolorbox}[colback=curs,colframe=curs,coltext=curstext,arc=2mm,boxrule=0pt,left=3mm,right=3mm,top=1.2mm,bottom=1.2mm]
  {\large\bfseries #1\quad #2}
  \end{tcolorbox}\vspace{-2pt}}

% ---------- estrelles de dificultat del color del curs (a Batxillerat no se'n posen) ----------
\@ifundefined{estrelles}{}{%
  \renewcommand{\estrelles}[1]{{\color{curs}\ifcase#1\or\faStar\faStar[regular]\faStar[regular]\or\faStar\faStar\faStar[regular]\or\faStar\faStar\faStar\fi}}}

% ---------- capçalera i peu de les pàgines ----------
% primera pàgina: logos institucionals al peu
\fancypagestyle{primera}{\fancyhf{}\renewcommand{\headrulewidth}{0pt}
  \fancyfoot[L]{\raisebox{-0.32\height}{\includegraphics[height=0.78cm]{logo-gva.png}}\hspace{5mm}%
    \raisebox{-0.32\height}{\includegraphics[height=0.78cm]{logo-vermellar.png}}}
  \fancyfoot[R]{\scriptsize\color{muted}\href{https://creativecommons.org/licenses/by-nc-sa/4.0/deed.ca}{\ccbyncsa}\ \ CC BY-NC-SA 4.0 · fentmates.org\qquad\small\thepage}}
% la resta: el símbol a l'esquerra i el curs a la dreta, amb el títol al centre com abans
\fancyhead[L]{\raisebox{-0.3\height}{\includegraphics[height=0.62cm]{\fm@marca}}}
\fancyhead[R]{\small\bfseries\color{cursfosc}\fm@cursnom}
\renewcommand{\headrule}{\color{curs!45}\hrule height \headrulewidth}
\makeatother

%   \fmcurs{1eso}
% i, al començament del document:
%   \capcalera{Matemàtiques}{1r ESO}{Fitxa d'exercicis · Unitat 2}{Divisibilitat}
% Opcional abans de \input: \def\fmneed{14} (espai mínim abans de cada secció).
% =====================================================================
\makeatletter
\def\fm@logodir{../../logo/pdf}
\providecommand{\fmneed}{10}
\def\fm@cursnom{}

% clau del curs, color i color del text damunt del color
\newcommand{\fm@curs}[3]{\@namedef{fm@color@#1}{#2}\@namedef{fm@text@#1}{#3}}
\fm@curs{1eso}{E6B84A}{navy}
\fm@curs{2eso}{3B82F6}{white}
\fm@curs{3eso}{FF7A50}{white}
\fm@curs{4eso-a}{A78BFA}{white}
\fm@curs{4eso-b}{6D28D9}{white}
\fm@curs{1batx-mi}{9B1C1C}{white}
\fm@curs{1batx-ccss}{EF6B6B}{white}
\fm@curs{2batx-mii}{1E6B47}{white}
\fm@curs{2batx-ccss}{2BB3A3}{white}

\newcommand{\fmcurs}[1]{%
  \@ifundefined{fm@color@#1}{\PackageError{fentmates}{Curs desconegut: #1}{}}{}%
  \edef\fm@tmp{\noexpand\definecolor{curs}{HTML}{\@nameuse{fm@color@#1}}%
               \noexpand\colorlet{curstext}{\@nameuse{fm@text@#1}}}\fm@tmp
  \colorlet{cursfosc}{curs!72!black}%
  \def\fm@logo{\fm@logodir/vertical-#1.pdf}%
  \def\fm@marca{\fm@logodir/marca-#1.pdf}}

% ---------- capçalera de la primera pàgina ----------
% #1 matèria (petit), #2 curs, #3 tipus de document i unitat, #4 títol
\newcommand{\capcalera}[4]{%
  \gdef\fm@cursnom{#2}%
  \par\noindent
  \begin{tcolorbox}[colback=curs,colframe=curs,boxrule=0pt,arc=2.2mm,width=3.5cm,height=2.15cm,
      valign=center,left=3mm,right=2mm,top=1mm,bottom=1mm,nobeforeafter,box align=center]
    \color{curstext}{\footnotesize #1}\par\vspace{1pt}{\fontsize{19}{21}\selectfont\bfseries #2\par}%
  \end{tcolorbox}\hspace{5mm}%
  \begin{minipage}[c]{\dimexpr\linewidth-3.5cm-5mm-2.9cm\relax}
    {\color{cursfosc}\small\bfseries\MakeUppercase{#3}}\par\vspace{2pt}
    {\color{navy}\fontsize{20}{23}\selectfont\bfseries #4\par}
  \end{minipage}\hfill
  \begin{minipage}[c]{2.6cm}\raggedleft\includegraphics[width=2.6cm]{\fm@logo}\end{minipage}\par
  \vspace{2pt}}

% ---------- barres de secció del color del curs ----------
\renewcommand{\seccio}[2]{\par\needspace{\fmneed\baselineskip}\bigskip
  \begin{tcolorbox}[colback=curs,colframe=curs,coltext=curstext,arc=2mm,boxrule=0pt,left=3mm,right=3mm,top=1.2mm,bottom=1.2mm]
  {\large\bfseries #1\quad #2}
  \end{tcolorbox}\vspace{-2pt}}

% ---------- estrelles de dificultat del color del curs (a Batxillerat no se'n posen) ----------
\@ifundefined{estrelles}{}{%
  \renewcommand{\estrelles}[1]{{\color{curs}\ifcase#1\or\faStar\faStar[regular]\faStar[regular]\or\faStar\faStar\faStar[regular]\or\faStar\faStar\faStar\fi}}}

% ---------- capçalera i peu de les pàgines ----------
% primera pàgina: logos institucionals al peu
\fancypagestyle{primera}{\fancyhf{}\renewcommand{\headrulewidth}{0pt}
  \fancyfoot[L]{\raisebox{-0.32\height}{\includegraphics[height=0.78cm]{logo-gva.png}}\hspace{5mm}%
    \raisebox{-0.32\height}{\includegraphics[height=0.78cm]{logo-vermellar.png}}}
  \fancyfoot[R]{\scriptsize\color{muted}\href{https://creativecommons.org/licenses/by-nc-sa/4.0/deed.ca}{\ccbyncsa}\ \ CC BY-NC-SA 4.0 · fentmates.org\qquad\small\thepage}}
% la resta: el símbol a l'esquerra i el curs a la dreta, amb el títol al centre com abans
\fancyhead[L]{\raisebox{-0.3\height}{\includegraphics[height=0.62cm]{\fm@marca}}}
\fancyhead[R]{\small\bfseries\color{cursfosc}\fm@cursnom}
\renewcommand{\headrule}{\color{curs!45}\hrule height \headrulewidth}
\makeatother

%   \fmcurs{1eso}
% i, al començament del document:
%   \capcalera{Matemàtiques}{1r ESO}{Fitxa d'exercicis · Unitat 2}{Divisibilitat}
% Opcional abans de \input: \def\fmneed{14} (espai mínim abans de cada secció).
% =====================================================================
\makeatletter
\def\fm@logodir{../../logo/pdf}
\providecommand{\fmneed}{10}
\def\fm@cursnom{}

% clau del curs, color i color del text damunt del color
\newcommand{\fm@curs}[3]{\@namedef{fm@color@#1}{#2}\@namedef{fm@text@#1}{#3}}
\fm@curs{1eso}{E6B84A}{navy}
\fm@curs{2eso}{3B82F6}{white}
\fm@curs{3eso}{FF7A50}{white}
\fm@curs{4eso-a}{A78BFA}{white}
\fm@curs{4eso-b}{6D28D9}{white}
\fm@curs{1batx-mi}{9B1C1C}{white}
\fm@curs{1batx-ccss}{EF6B6B}{white}
\fm@curs{2batx-mii}{1E6B47}{white}
\fm@curs{2batx-ccss}{2BB3A3}{white}

\newcommand{\fmcurs}[1]{%
  \@ifundefined{fm@color@#1}{\PackageError{fentmates}{Curs desconegut: #1}{}}{}%
  \edef\fm@tmp{\noexpand\definecolor{curs}{HTML}{\@nameuse{fm@color@#1}}%
               \noexpand\colorlet{curstext}{\@nameuse{fm@text@#1}}}\fm@tmp
  \colorlet{cursfosc}{curs!72!black}%
  \def\fm@logo{\fm@logodir/vertical-#1.pdf}%
  \def\fm@marca{\fm@logodir/marca-#1.pdf}}

% ---------- capçalera de la primera pàgina ----------
% #1 matèria (petit), #2 curs, #3 tipus de document i unitat, #4 títol
\newcommand{\capcalera}[4]{%
  \gdef\fm@cursnom{#2}%
  \par\noindent
  \begin{tcolorbox}[colback=curs,colframe=curs,boxrule=0pt,arc=2.2mm,width=3.5cm,height=2.15cm,
      valign=center,left=3mm,right=2mm,top=1mm,bottom=1mm,nobeforeafter,box align=center]
    \color{curstext}{\footnotesize #1}\par\vspace{1pt}{\fontsize{19}{21}\selectfont\bfseries #2\par}%
  \end{tcolorbox}\hspace{5mm}%
  \begin{minipage}[c]{\dimexpr\linewidth-3.5cm-5mm-2.9cm\relax}
    {\color{cursfosc}\small\bfseries\MakeUppercase{#3}}\par\vspace{2pt}
    {\color{navy}\fontsize{20}{23}\selectfont\bfseries #4\par}
  \end{minipage}\hfill
  \begin{minipage}[c]{2.6cm}\raggedleft\includegraphics[width=2.6cm]{\fm@logo}\end{minipage}\par
  \vspace{2pt}}

% ---------- barres de secció del color del curs ----------
\renewcommand{\seccio}[2]{\par\needspace{\fmneed\baselineskip}\bigskip
  \begin{tcolorbox}[colback=curs,colframe=curs,coltext=curstext,arc=2mm,boxrule=0pt,left=3mm,right=3mm,top=1.2mm,bottom=1.2mm]
  {\large\bfseries #1\quad #2}
  \end{tcolorbox}\vspace{-2pt}}

% ---------- estrelles de dificultat del color del curs (a Batxillerat no se'n posen) ----------
\@ifundefined{estrelles}{}{%
  \renewcommand{\estrelles}[1]{{\color{curs}\ifcase#1\or\faStar\faStar[regular]\faStar[regular]\or\faStar\faStar\faStar[regular]\or\faStar\faStar\faStar\fi}}}

% ---------- capçalera i peu de les pàgines ----------
% primera pàgina: logos institucionals al peu
\fancypagestyle{primera}{\fancyhf{}\renewcommand{\headrulewidth}{0pt}
  \fancyfoot[L]{\raisebox{-0.32\height}{\includegraphics[height=0.78cm]{logo-gva.png}}\hspace{5mm}%
    \raisebox{-0.32\height}{\includegraphics[height=0.78cm]{logo-vermellar.png}}}
  \fancyfoot[R]{\scriptsize\color{muted}\href{https://creativecommons.org/licenses/by-nc-sa/4.0/deed.ca}{\ccbyncsa}\ \ CC BY-NC-SA 4.0 · fentmates.org\qquad\small\thepage}}
% la resta: el símbol a l'esquerra i el curs a la dreta, amb el títol al centre com abans
\fancyhead[L]{\raisebox{-0.3\height}{\includegraphics[height=0.62cm]{\fm@marca}}}
\fancyhead[R]{\small\bfseries\color{cursfosc}\fm@cursnom}
\renewcommand{\headrule}{\color{curs!45}\hrule height \headrulewidth}
\makeatother

\fmcurs{1eso}

\begin{document}
\thispagestyle{primera}

\capcalera{Matemàtiques}{1r ESO}{Solucions de la fitxa · Unitat 8}{Àrees i perímetres}
\vspace{8pt}

\begin{tcolorbox}[colback=soft,colframe=soft,arc=2mm,left=3mm,right=3mm,top=1.5mm,bottom=1.5mm]
\small Intenta resoldre cada exercici abans de mirar-ne la solució. Si no t'ix el mateix resultat, repassa els passos o consulta els apunts interactius de la unitat: hi trobaràs eines per a canviar d'unitats i per a explorar les àrees i els perímetres de totes les figures.
\end{tcolorbox}

% =====================================================
\seccio{1}{Perímetre i àrea. Unitats de mesura}

\ex{1} \begin{tasks}(3)
  \task $A=\res{6\text{ cm}^2}$, $P=\res{10\text{ cm}}$
  \task $A=\res{7\text{ cm}^2}$, $P=\res{16\text{ cm}}$
  \task $A=\res{5\text{ cm}^2}$, $P=\res{12\text{ cm}}$
\end{tasks}

\ex{1} \begin{tasks}(5)
  \task $\res{4700\text{ m}}$
  \task $\res{3,8\text{ m}}$
  \task $\res{2,5\text{ m}}$
  \task $\res{900\text{ m}}$
  \task $\res{1,25\text{ m}}$
\end{tasks}

\ex{1} \begin{tasks}(4)
  \task $\res{320\text{ cm}}$
  \task $\res{4,5\text{ cm}}$
  \task $\res{80\text{ dam}}$
  \task $\res{0,76\text{ hm}}$
\end{tasks}

\ex{1} \begin{tasks}(3)
  \task $\res{500\text{ cm}^2}$
  \task $\res{240\text{ dm}^2}$
  \task $\res{3,5\text{ m}^2}$
  \task $\res{7000\text{ m}^2}$
  \task $\res{800\text{ hm}^2}$
\end{tasks}

\ex{2} \begin{tasks}(4)
  \task $\res{35\,000\text{ m}^2}$
  \task $\res{0,45\text{ ha}}$
  \task $\res{1200\text{ m}^2}$
  \task $\res{25\text{ ha}}$
\end{tasks}

\ex{2} Ho passem tot a m$^2$: $50\,000$, $40\,000$, $52\,000$ i $48\,000$. \ \rest{$4\text{ ha}<48\,000\text{ m}^2<0,05\text{ km}^2<520\text{ dam}^2$}

\ex{2} $2,4\text{ ha}=24\,000\text{ m}^2$; \ $24\,000\cdot18=\res{432\,000\text{ €}}$

\ex{3} La base i l'altura sumen 10: \ $1\times9$ ($A=9$), $2\times8$ ($A=16$), $3\times7$ ($A=21$), $4\times6$ ($A=24$) i $5\times5$ ($A=25$). L'àrea més gran és la del \rest{quadrat de 5\,cm}: $25\text{ cm}^2$.


% =====================================================
\seccio{2}{Perímetre dels polígons}

\ex{1} \begin{tasks}(2)
  \task $\res{27,5\text{ cm}}$
  \task $4\cdot8,5=\res{34\text{ cm}}$
  \task $2\cdot15+2\cdot9=\res{48\text{ m}}$
  \task $4\cdot6,3=\res{25,2\text{ cm}}$
\end{tasks}

\ex{1} \begin{tasks}(2)
  \task $5\cdot7=\res{35\text{ cm}}$
  \task $6\cdot4,5=\res{27\text{ cm}}$
  \task $8\cdot3,2=\res{25,6\text{ cm}}$
  \task $10\cdot12=\res{120\text{ mm}}$
\end{tasks}

\ex{1} $P=5+3+4+3,5+3=\res{18,5\text{ cm}}$

\ex{2} \begin{tasks}(3)
  \task $46:4=\res{11,5\text{ cm}}$
  \task $51:6=\res{8,5\text{ cm}}$
  \task $13,5:3=\res{4,5\text{ m}}$
\end{tasks}

\ex{2} Les dues bases fan $40$\,cm; queden $64-40=24$\,cm per a les dues altures: \ $h=24:2=\res{12\text{ cm}}$.

\ex{2} $P=2\cdot120+2\cdot85=410$\,m. \ Amb tres fils: $3\cdot410=\res{1230\text{ m}}$. \ Preu: $1230\cdot0,35=\res{430,50\text{ €}}$.

\ex{3} El perímetre equival a $3+1+3+1=8$ altures: $h=96:8=\res{12\text{ cm}}$ i $b=3\cdot12=\res{36\text{ cm}}$.

\ex{3} El perímetre equival a $1+2+2=5$ bases: base $=40:5=\res{8\text{ cm}}$ i els costats iguals, \res{16\text{ cm}} cadascun.


% =====================================================
\seccio{3}{Àrea del quadrat i del rectangle}

\ex{1} \begin{tasks}(2)
  \task $A=\res{144\text{ cm}^2}$, $P=\res{48\text{ cm}}$
  \task $A=\res{120\text{ cm}^2}$, $P=\res{46\text{ cm}}$
  \task $A=\res{30\text{ m}^2}$, $P=\res{23\text{ m}}$
  \task $A=\res{6,25\text{ m}^2}$, $P=\res{10\text{ m}}$
\end{tasks}

\ex{1} $h=84:12=\res{7\text{ cm}}$; \ $P=2\cdot12+2\cdot7=\res{38\text{ cm}}$

\ex{1} \begin{tasks}(1)
  \task $11^2=121\Rightarrow l=\res{11\text{ m}}$; \ $P=\res{44\text{ m}}$
  \task $l=28:4=7$\,cm; \ $A=7^2=\res{49\text{ cm}^2}$
\end{tasks}

\ex{2} $A=5,4\cdot3,5=18,9\text{ m}^2$; \ $18,9\cdot24=\res{453,60\text{ €}}$

\ex{2} $6\cdot2,5-1,5\cdot1,2=15-1,8=13,2\text{ m}^2$. \ $13,2:5=2,64$: cal comprar \rest{3 pots}.

\ex{2} Caben $800:50=16$ rajoles al llarg i $600:50=12$ a l'ample: \ $16\cdot12=\res{192}$ rajoles.

\ex{3} $P$: de 20\,cm a 40\,cm, \rest{es duplica}. \ $A$: de 25\,cm$^2$ a 100\,cm$^2$, \rest{es multiplica per 4}.

\ex{3} $1\times36$ ($P=74$), $2\times18$ ($P=40$), $3\times12$ ($P=30$), $4\times9$ ($P=26$), $6\times6$ ($P=24$). El de perímetre més menut és el \rest{quadrat de 6\,cm}.


% =====================================================
\seccio{4}{Àrea del romboide i del triangle}

\ex{1} \begin{tasks}(2)
  \task $12\cdot7=\res{84\text{ cm}^2}$
  \task $8,5\cdot4=\res{34\text{ m}^2}$
\end{tasks}

\ex{1} \begin{tasks}(3)
  \task $\frac{14\cdot9}{2}=\res{63\text{ cm}^2}$
  \task $\frac{7\cdot5}{2}=\res{17,5\text{ cm}^2}$
  \task $\frac{6,4\cdot3,5}{2}=\res{11,2\text{ m}^2}$
\end{tasks}

\ex{1} \begin{tasks}(1)
  \task $\frac{6\cdot4}{2}=\res{12\text{ cm}^2}$
  \task $\frac{8\cdot5}{2}=\res{20\text{ cm}^2}$
  \task $\frac{4\cdot5}{2}=\res{10\text{ cm}^2}$ (l'altura cau fora)
\end{tasks}

\ex{2} Els catets fan de base i altura: \ $A=\frac{9\cdot12}{2}=\res{54\text{ cm}^2}$. \ $P=9+12+15=\res{36\text{ cm}}$.

\ex{2} \begin{tasks}(1)
  \task $h=96:12=\res{8\text{ cm}}$
  \task $b\cdot9=90\Rightarrow b=\res{10\text{ m}}$
\end{tasks}

\ex{2} Base 8\,cm (sobre la línia de baix) i altura 5\,cm: \ $A=\frac{8\cdot5}{2}=\res{20\text{ cm}^2}$.

\ex{3} $A=\frac{14\cdot12}{2}=\res{84\text{ cm}^2}$. \ Amb base 15: $\frac{15\cdot h}{2}=84\Rightarrow 15\cdot h=168\Rightarrow h=\res{11,2\text{ cm}}$.

\ex{3} Tots tenen la mateixa base (10\,cm) i la mateixa altura (6\,cm, la distància entre les paral·leles): $A=\frac{10\cdot6}{2}=\res{30\text{ cm}^2}$ \rest{tots tres}.


% =====================================================
\seccio{5}{Àrea del rombe i del trapezi}

\ex{1} \begin{tasks}(3)
  \task $\frac{16\cdot10}{2}=\res{80\text{ cm}^2}$
  \task $\frac{9\cdot6}{2}=\res{27\text{ m}^2}$
  \task $\frac{7,5\cdot4}{2}=\res{15\text{ cm}^2}$
\end{tasks}

\ex{1} \begin{tasks}(3)
  \task $\frac{(12+8)\cdot5}{2}=\res{50\text{ cm}^2}$
  \task $\frac{(9+4)\cdot6}{2}=\res{39\text{ m}^2}$
  \task $\frac{17\cdot4}{2}=\res{34\text{ cm}^2}$
\end{tasks}

\ex{2} $l=52:4=\res{13\text{ cm}}$; \ $A=\frac{24\cdot10}{2}=\res{120\text{ cm}^2}$

\ex{2} $A=\frac{(18+10)\cdot3}{2}=\res{42\text{ cm}^2}$; \ $P=18+10+5+5=\res{38\text{ cm}}$

\ex{2} \begin{tasks}(1)
  \task $15\cdot d=120\Rightarrow d=\res{8\text{ cm}}$
  \task $(14+10)\cdot h=144\Rightarrow 24\cdot h=144\Rightarrow h=\res{6\text{ m}}$
\end{tasks}

\ex{3} $A=\frac{1,2\cdot0,8}{2}=\res{0,48\text{ m}^2}=\res{4800\text{ cm}^2}$

\ex{3} $A=\frac{(85+55)\cdot40}{2}=2800\text{ m}^2=\res{0,28\text{ ha}}$. \ Preu: $2800\cdot12=\res{33\,600\text{ €}}$.


% =====================================================
\seccio{6}{Àrea dels polígons regulars i dels polígons qualssevol}

\ex{1} $P=5\cdot10=\res{50\text{ cm}}$; \ $A=\frac{50\cdot6,88}{2}=\res{172\text{ cm}^2}$

\ex{1} $P=6\cdot8=\res{48\text{ cm}}$; \ $A=\frac{48\cdot6,93}{2}=\res{166,32\text{ cm}^2}$

\ex{2} $l=40:8=\res{5\text{ cm}}$; \ $A=\frac{40\cdot6,04}{2}=\res{120,8\text{ cm}^2}$

\ex{2} $\frac{36\cdot a}{2}=93,6\Rightarrow 36\cdot a=187,2\Rightarrow a=\res{5,2\text{ cm}}$

\ex{2} Rectangle $6\times4$ i triangle de base 6 i altura 2: \ $24+\frac{6\cdot2}{2}=24+6=\res{30\text{ cm}^2}$

\ex{3} Rectangle $7\times6=42$. Triangles que sobren: $\frac{7\cdot2}{2}=7$, \ $\frac{4\cdot4}{2}=8$ \ i \ $\frac{3\cdot6}{2}=9$. \ $A=42-7-8-9=\res{18\text{ cm}^2}$

\ex{3} Rajola: $A=\frac{60\cdot8,66}{2}=259,8\text{ cm}^2$. \ $25\text{ m}^2=250\,000\text{ cm}^2$; \ $250\,000:259,8\approx962,3$: calen \rest{963 rajoles}.


% =====================================================
\seccio{7}{La circumferència i la seua longitud}

\ex{1} \begin{tasks}(4)
  \task $a$: \rest{radi}
  \task $b$: \rest{diàmetre}
  \task $c$: \rest{corda}
  \task $d$: \rest{arc}
\end{tasks}

\ex{1} \begin{tasks}(2)
  \task $2\cdot3,14\cdot4=\res{25,12\text{ cm}}$
  \task $2\cdot3,14\cdot9=\res{56,52\text{ m}}$
  \task $3,14\cdot20=\res{62,8\text{ cm}}$
  \task $3,14\cdot7=\res{21,98\text{ mm}}$
\end{tasks}

\ex{2} \begin{tasks}(1)
  \task $d=31,4:3,14=10$\,cm; \ $r=\res{5\text{ cm}}$
  \task $d=75,36:3,14=\res{24\text{ m}}$; \ $r=\res{12\text{ m}}$
\end{tasks}

\ex{2} $L=3,14\cdot70=\res{219,8\text{ cm}}$. \ $1\text{ km}=100\,000\text{ cm}$; \ $100\,000:219,8\approx455$ \ \rest{unes 455 voltes}.

\ex{2} \begin{tasks}(3)
  \task $\frac{2\cdot3,14\cdot8\cdot90}{360}=\res{12,56\text{ cm}}$
  \task $\frac{2\cdot3,14\cdot12\cdot60}{360}=\res{12,56\text{ cm}}$
  \task $\frac{2\cdot3,14\cdot6\cdot150}{360}=\res{15,7\text{ cm}}$
\end{tasks}

\ex{3} \begin{tasks}(2)
  \task $\frac{3,14\cdot10}{2}+10=15,7+10=\res{25,7\text{ cm}}$
  \task $\frac{2\cdot3,14\cdot4}{4}+4+4=6,28+8=\res{14,28\text{ cm}}$
\end{tasks}

\ex{3} En 20 minuts gira $120^\circ$: \ $\frac{2\cdot3,14\cdot15\cdot120}{360}=\res{31,4\text{ cm}}$. \ En una hora: $2\cdot3,14\cdot15=\res{94,2\text{ cm}}$.


% =====================================================
\seccio{8}{L'àrea del cercle i de les figures circulars}

\ex{1} \begin{tasks}(2)
  \task $3,14\cdot36=\res{113,04\text{ cm}^2}$
  \task $3,14\cdot6,25=\res{19,625\text{ m}^2}$
  \task $3,14\cdot7^2=\res{153,86\text{ cm}^2}$
  \task $3,14\cdot15^2=\res{706,5\text{ cm}^2}$
\end{tasks}

\ex{2} \begin{tasks}(2)
  \task $\frac{3,14\cdot100}{2}=\res{157\text{ cm}^2}$
  \task $\frac{3,14\cdot36}{4}=\res{28,26\text{ cm}^2}$
\end{tasks}

\ex{2} \begin{tasks}(3)
  \task $\frac{3,14\cdot64\cdot45}{360}=\res{25,12\text{ cm}^2}$
  \task $\frac{3,14\cdot81\cdot120}{360}=\res{84,78\text{ cm}^2}$
  \task $\frac{3,14\cdot16\cdot270}{360}=\res{37,68\text{ cm}^2}$
\end{tasks}

\ex{2} \begin{tasks}(2)
  \task $3,14\cdot(49-16)=3,14\cdot33=\res{103,62\text{ cm}^2}$
  \task $3,14\cdot(100-36)=3,14\cdot64=\res{200,96\text{ m}^2}$
\end{tasks}

\ex{2} \begin{tasks}(1)
  \task $r^2=78,5:3,14=25\Rightarrow r=\res{5\text{ cm}}$
  \task $r=18,84:6,28=3$\,cm; \ $A=3,14\cdot9=\res{28,26\text{ cm}^2}$
\end{tasks}

\ex{3} Corona de radis $R=6$\,cm i $r=0,75$\,cm: \ $3,14\cdot(36-0,5625)\approx\res{111,27\text{ cm}^2}$

\ex{3} Gran: $3,14\cdot16^2=803,84\text{ cm}^2$. \ Dues menudes: $2\cdot3,14\cdot11^2=759,88\text{ cm}^2$. \ \rest{La pizza gran.}


% =====================================================
\seccio{9}{Figures compostes}

\ex{1} Completem el rectangle $10\times7$ i restem el tros $4\times3$: \ $A=70-12=\res{58\text{ cm}^2}$. \ $P=10+4+4+3+6+7=\res{34\text{ cm}}$.

\ex{1} $A=8\cdot5+\frac{8\cdot3}{2}=40+12=\res{52\text{ m}^2}$

\ex{2} Semicercle de radi 2\,cm: \ $A=6\cdot4+\frac{3,14\cdot2^2}{2}=24+6,28=\res{30,28\text{ cm}^2}$. \ $P=6+6+4+\frac{2\cdot3,14\cdot2}{2}=16+6,28=\res{22,28\text{ cm}}$

\ex{2} Els quatre quarts formen un cercle sencer de radi 4\,cm: \ $A=8^2-3,14\cdot4^2=64-50,24=\res{13,76\text{ cm}^2}$

\ex{2} Rectangle exterior: $(10+3)\times(6+3)=13\cdot9=117\text{ m}^2$. \ $A=117-60=\res{57\text{ m}^2}$

\ex{2} Els dos semicercles fan un cercle de radi 25\,m: \ $A=80\cdot50+3,14\cdot25^2=4000+1962,5=\res{5962,5\text{ m}^2}$. \ $P=2\cdot80+3,14\cdot50=160+157=\res{317\text{ m}}$

\ex{3} Centre: $3,14\cdot9=\res{28,26\text{ cm}^2}$. \ Zona del mig: $3,14\cdot(36-9)=\res{84,78\text{ cm}^2}$. \ Zona de fora: $3,14\cdot(81-36)=\res{141,3\text{ cm}^2}$.

\ex{3} La diagonal divideix la fulla en dues meitats. Cada meitat és un quart de cercle menys un triangle: $\frac{3,14\cdot36}{4}-\frac{6\cdot6}{2}=28,26-18=10,26$. \ $A=2\cdot10,26=\res{20,52\text{ cm}^2}$

\ex{3} $A=1,2\cdot1,5+\frac{3,14\cdot0,6^2}{2}=1,8+0,5652=2,3652\text{ m}^2$. \ $2,3652\cdot45\approx\res{106,43\text{ €}}$


% =====================================================
\seccio{10}{El teorema de Pitàgores}

\ex{1} \begin{tasks}(1)
  \task Catets $x$ i $y$; hipotenusa $z$: \ $\res{z^2=x^2+y^2}$
  \task Catets $n$ i $p$; hipotenusa $m$: \ $\res{m^2=n^2+p^2}$
  \task Catets $r$ i $s$; hipotenusa $t$: \ $\res{t^2=r^2+s^2}$
\end{tasks}

\ex{1} \begin{tasks}(2)
  \task $\sqrt{9+16}=\sqrt{25}=\res{5\text{ cm}}$
  \task $\sqrt{36+64}=\sqrt{100}=\res{10\text{ m}}$
  \task $\sqrt{25+144}=\sqrt{169}=\res{13\text{ cm}}$
  \task $\sqrt{81+144}=\sqrt{225}=\res{15\text{ dm}}$
\end{tasks}

\ex{1} \begin{tasks}(2)
  \task $\sqrt{100-36}=\sqrt{64}=\res{8\text{ cm}}$
  \task $\sqrt{169-144}=\sqrt{25}=\res{5\text{ m}}$
  \task $\sqrt{289-225}=\sqrt{64}=\res{8\text{ cm}}$
  \task $\sqrt{625-49}=\sqrt{576}=\res{24\text{ cm}}$
\end{tasks}

\ex{2} \begin{tasks}(2)
  \task $c=\sqrt{4+25}=\sqrt{29}\approx\res{5,39\text{ cm}}$
  \task $c=\sqrt{16+36}=\sqrt{52}\approx\res{7,21\text{ cm}}$
  \task $b=\sqrt{64-9}=\sqrt{55}\approx\res{7,42\text{ cm}}$
  \task $b=\sqrt{100-49}=\sqrt{51}\approx\res{7,14\text{ cm}}$
\end{tasks}

\ex{2} \begin{tasks}(2)
  \task $10^2=100$; \ $36+64=100$: \ \rest{sí}
  \task $9^2=81$; \ $25+49=74$: \ \rest{no}
  \task $17^2=289$; \ $64+225=289$: \ \rest{sí}
  \task $15^2=225$; \ $100+121=221$: \ \rest{no}
\end{tasks}

\ex{2} L'escala és la hipotenusa: \ $h=\sqrt{5^2-1,4^2}=\sqrt{25-1,96}=\sqrt{23,04}=\res{4,8\text{ m}}$

\ex{2} Els dos trajectes són els catets: \ $d=\sqrt{12^2+9^2}=\sqrt{144+81}=\sqrt{225}=\res{15\text{ km}}$

\ex{3} Triangle rectangle de catets 6\,cm i 8\,cm: \ $AB=\sqrt{36+64}=\sqrt{100}=\res{10\text{ cm}}$

\ex{3} Cada cable: $\sqrt{12^2+5^2}=\sqrt{169}=13$\,m. \ Quatre cables: $4\cdot13=\res{52\text{ m}}$


% =====================================================
\seccio{11}{Pitàgores per a calcular àrees i perímetres}

\ex{1} \begin{tasks}(3)
  \task $\sqrt{64+36}=\res{10\text{ cm}}$
  \task $\sqrt{576+49}=\res{25\text{ m}}$
  \task $\sqrt{16+16}=\sqrt{32}\approx\res{5,66\text{ cm}}$
\end{tasks}

\ex{1} L'altura parteix la base en dues meitats de 8\,cm: \ $h=\sqrt{10^2-8^2}=\sqrt{36}=\res{6\text{ cm}}$. \ $A=\frac{16\cdot6}{2}=\res{48\text{ cm}^2}$. \ $P=16+10+10=\res{36\text{ cm}}$

\ex{2} Les semidiagonals fan 9 i 12\,cm: \ $l=\sqrt{81+144}=\sqrt{225}=\res{15\text{ cm}}$. \ $P=4\cdot15=\res{60\text{ cm}}$. \ $A=\frac{18\cdot24}{2}=\res{216\text{ cm}^2}$

\ex{2} Triangle de catets $12-7=5$ i 12: \ $c=\sqrt{25+144}=\sqrt{169}=13$\,cm. \ $P=12+7+12+13=\res{44\text{ cm}}$. \ $A=\frac{(12+7)\cdot12}{2}=\res{114\text{ cm}^2}$

\ex{2} A cada costat sobren $(16-4):2=6$\,cm: \ $h=\sqrt{100-36}=\res{8\text{ cm}}$. \ $A=\frac{(16+4)\cdot8}{2}=\res{80\text{ cm}^2}$. \ $P=16+4+10+10=\res{40\text{ cm}}$

\ex{2} El radi és igual al costat (4\,cm): \ $a=\sqrt{4^2-2^2}=\sqrt{12}\approx\res{3,46\text{ cm}}$. \ $P=6\cdot4=24$\,cm; \ $A=\frac{24\cdot3,46}{2}=\res{41,52\text{ cm}^2}$

\ex{3} Costat: $18:3=6$\,cm. \ $h=\sqrt{6^2-3^2}=\sqrt{27}\approx\res{5,20\text{ cm}}$. \ $A=\frac{6\cdot5,2}{2}=\res{15,6\text{ cm}^2}$

\ex{3} La diagonal del quadrat és un diàmetre (8\,cm): \ $l^2+l^2=64\Rightarrow l^2=32$, que és l'àrea del quadrat. \ $A=3,14\cdot4^2-32=50,24-32=\res{18,24\text{ cm}^2}$

\ex{3} Cada costat de la teulada: $\sqrt{4^2+3^2}=\sqrt{25}=5$\,m. \ $P=8+5+5+5+5=\res{28\text{ m}}$. \ $A=8\cdot5+\frac{8\cdot3}{2}=\res{52\text{ m}^2}$


\end{document}
